\documentclass[aspectratio=169]{beamer}
\usetheme{Madrid}
\usecolortheme{default}

\usepackage{amsmath,amssymb}
\usepackage{tikz}
\usetikzlibrary{arrows.meta, positioning}

\providecommand{\labelenumi}{}
\providecommand{\labelenumii}{}

\input{simple-type-theory-def.tex}
\input{alonzo-notation.tex}

\newcommand{\nil}{\mName{nil}}
\newcommand{\cons}{\mName{cons}}
\newcommand{\some}{\mName{some}}
\newcommand{\none}{\mName{none}}
\newcommand{\leaf}{\mName{leaf}}
\newcommand{\node}{\mName{node}}
\newcommand{\true}{\mName{true}}
\newcommand{\false}{\mName{false}}

\title{OCaml Core Types in Alonzo}
\subtitle{CAS 760 --- Development 2 Project}
\author{Suleyman Kiani \and Hamza}
\institute{McMaster University}
\date{April 16, 2026}

\begin{document}

\begin{frame}
\titlepage
\end{frame}

\begin{frame}{Outline}
\tableofcontents
\end{frame}

%% Presentation outline / guideline
%%
%% Planned structure (~15 min total):
%%
%% 1. Motivation (~2 min)
%%    - Why OCaml types? Connection between PL types and inductive types
%%    - Each OCaml algebraic data type maps to an inductive type (Ch 10)
%%    - Gives nontrivial complexity: 4 theories + 2 developments + 1 morphism
%%
%% 2. Theory graph overview (~1 min)
%%    - Show the full diagram, explain the module relationships
%%
%% 3. Walk through each module (~8 min)
%%    - BOOL-TYPE: simplest case, enumeration type (Ex 10.5)
%%    - OPTION-TYPE: sum type with one constructor carrying data
%%    - LIST-TYPE: recursive type, parallels PA (nil<->0, cons<->S)
%%    - LIST-DEV: derived ops (length, append, reverse), show one proof
%%    - TREE-TYPE: binary recursive type (Ex 10.7)
%%    - TREE-DEV: derived ops (size, height, mirror)
%%
%% 4. Theory morphism (~2 min)
%%    - PA -> LIST-TYPE, explain the mapping and what Thm 14.16 gives us
%%
%% 5. Conclusion + Questions (~2 min)
%%
%% This covers the scope because:
%% - Multiple theory defs (BOOL, OPTION, LIST, TREE)
%% - Development defs with Defs + Thms (LIST-DEV, TREE-DEV)
%% - A theory morphism (PA -> LIST)
%% - Nontrivial complexity (7 modules total)
%% - Follows the monoids paper structure
%%
%% Lessons from Project 1 to address in the talk:
%% - Show that derived ops are in developments (not as ext constants)
%% - Show Def validations in proofs
%% - Emphasize correct use of TOTAL thms

\section{Motivation}
\section{Theory Graph}
\section{Boolean Type}
\section{Option Type}
\section{List Type and Development}
\section{Binary Tree Type and Development}
\section{Theory Morphism}
\section{Conclusion}

\begin{frame}{Questions?}
\centering
\Large Thank you!
\end{frame}

\end{document}
